\documentclass[../main]{subfiles}
\begin{document}
\chapter{Motor paso a paso}
\normalsize
El motor paso a paso (\emph{Stepper} en inglés) conocido también como motor de pasos es un dispositivo electromecánico que convierte una serie de impulsos eléctricos en desplazamientos angulares discretos, lo que significa que es capaz de girar una cantidad de grados (paso o medio paso) dependiendo de sus entradas de control.

\video{https://www.youtube.com/watch?v=b_-PQCjyRRQ}{Funcionamiento de motor paso a paso.}

El motor paso a paso se comporta de la misma manera que un conversor digital-analógico (D/A) y puede ser gobernado por impulsos procedentes de sistemas digitales. Este motor presenta las ventajas de tener precisión y repetitividad en cuanto al posicionamiento. Entre sus principales aplicaciones destacan los robots, drones, radio control, impresoras digitales, automatización, foto componedoras, pre prensa, etc.

\img{images/stepper/stepper}{Motor paso a paso.}

\section{Tipos de motores paso a paso}
Existen 3 tipos fundamentales de motores paso a paso: el motor de reluctancia variable, el motor de magnetización permanente, y el motor híbrido.
\begin{itemize}
  \item El motor de pasos de reluctancia variable (VR): Tiene un rotor multipolar de hierro y un estator devanado, opcionalmente laminado. Rota cuando el (o los) diente (s) más cercano (s) del rotor es (o son) atraído (s) a la (s) bobina (s) del estator energizada (s) (obteniéndose por lo tanto, la ruta de menor reluctancia). La respuesta de este motor es muy rápida, pero la inercia permitida en la carga es pequeña. Cuando los devanados no están energizados, el par estático de este tipo de motor es cero.
  \item El motor de pasos de rotor de imán permanente: Permite mantener un par diferente de cero cuando el motor no está energizado. Dependiendo de la construcción del motor, es típicamente posible obtener pasos angulares de $7,5^\circ$, $11,25^\circ$, $15^\circ$, $18^\circ$, $45^\circ$ o $90^\circ$. El ángulo de rotación se determina por el número de polos en el estator.
  \item El motor de pasos híbrido: Se caracteriza por tener varios dientes en el estator y en el rotor, el rotor con un imán concéntrico magnetizado axialmente alrededor de su eje. Se puede ver que esta configuración es una mezcla de los tipos de reluctancia variable e imán permanente. Este tipo de motor tiene una alta precisión y alto par, se puede configurar para suministrar un paso angular tan pequeño como $1,8^\circ$.
\end{itemize}

\section{Motores paso a paso unipolares}
\img{images/stepper/unipolar}{Modelo conceptual de motor unipolar.}
Estos motores suelen tener 5 o 6 cables de salida dependiendo de su conexión interna. Este tipo se caracteriza por ser más simple de controlar, estos utilizan un cable común a la fuente de alimentación y posteriormente se van colocando las otras líneas a tierra en un orden específico para generar cada paso, si tienen 6 cables es porque cada par de bobinas tienen un común separado, si tiene 5 cables es porque las cuatro bobinas tienen un polo común; un motor unipolar de 6 cables puede ser usado como un motor bipolar si se deja las líneas del común al aire.

\section{Motores paso a paso bipolares}
Estos tienen generalmente 4 cables de salida. Necesitan ciertos trucos para ser controlados debido a que requieren del cambio de dirección de flujo de corriente a través de las bobinas en la secuencia apropiada para realizar un movimiento.

Secuencia de rotación en un motor Bipolar: Obsérvese cómo la variación de la dirección del campo magnético creado en el estator producirá movimiento de seguimiento de parte del rotor del imán permanente, el cual intentará alinearse con el campo magnético inducido por las bobinas que excitan los electroimanes (en este caso A y B). Vcc es la alimentación de corriente continua (por ejemplo: $\qty{5}{V}$, $\qty{12}{V}$, $\qty{24}{V}$ \dots)

\section{Control de las bobinas}
\img{images/stepper/puenteh}{Topología de puente en H para las bobinas A y B}
\img{images/stepper/puenteh2}{Variación de la alimentación según los transistores T1, T2, T3, T4}

Para el control del motor paso a paso de este tipo (bipolar), se establece el principio de Puente H, si se activan T1 y T4, permiten la alimentación en un sentido; si cambiamos el sentido de la alimentación activando T2 y T3, cambiaremos el sentido de alimentación y el sentido de la corriente.

\info{Velocidad de rotación}{
  La velocidad de rotación en R.P.M. (revoluciones por minuto) viene definida por la ecuación:
  \begin{equation}
    N = 60 \cdot \frac{f}{n}
  \end{equation}

  Donde:
  \begin{itemize}
    \item $f$: frecuencia del tren de impulsos \unit{Hz}
    \item $n$: Número de polos que forman el motor
  \end{itemize}

  \begin{enumerate}
    \item \textbf{Velocidad angular ($\omega$)}:
    \[ \omega = \frac{2\pi \cdot N}{60} \]

    \item \textbf{Resolución del paso ($\alpha$)}:
    \[ \alpha = \frac{360}{S} \]

    \item \textbf{Torque ($\tau$)}:
    \[ \tau = F \cdot r \]

    \item \textbf{Corriente (I)}:
    \[ I = \frac{V}{R} \]
  \end{enumerate}
}

Si bien hay que decir que para estos motores, la máxima frecuencia admisible suele estar alrededor de los \qty{625}{Hz}, en caso de que la frecuencia de pulsos sea demasiado elevada, el motor puede reaccionar en alguna de las siguientes maneras:
\begin{itemize}
  \item No realizar ningún movimiento en absoluto.
  \item Comenzar a vibrar pero sin llegar a girar.
  \item Girar erráticamente.
  \item Girar en sentido opuesto.
  \item Perder potencia
\end{itemize}

Como ayuda es recomendable que también se coloque a disposición un simulador o circuito para probar estos motores paso a paso para descartar fallas en ello.

\actividad{Investigar y Responder}
\begin{enumerate}
  \item ¿Cuáles son las ventajas y desventajas de los motores paso a paso en comparación con los motores de corriente continua?
  \item ¿Cuáles son las aplicaciones típicas de los motores paso a paso?
  \item ¿Cuáles son los diferentes tipos de motores paso a paso disponibles en el mercado?
  \item ¿Cuál es la composición básica de un motor paso a paso?
  \item ¿Qué significa que un motor paso a paso tenga reluctancia variable?
  \item ¿Cuáles son las principales diferencias entre los motores paso a paso bipolares y unipolares?
  \item ¿Cómo se realiza la identificación y conexión de las bobinas en motores paso a paso bipolares y unipolares?
  \item ¿Qué es un controlador de motor paso a paso y qué características eléctricas son importantes al elegir uno?
  \item ¿Cuáles son las ventajas del ULN2003 como controlador de motor paso a paso y cuántos motores puede controlar?
  \item ¿Cuál es la diferencia entre el ULN2003 y el L293D en términos de control de motores paso a paso?
\end{enumerate}

\actividad{Resolver}
\begin{enumerate}
  \item Un motor paso a paso tiene una velocidad angular de $\omega = \qty{200}{rad/s}$. ¿Cuál es su velocidad en RPM (revoluciones por minuto)?
  \item Si un motor paso a paso tiene una resolución del paso de $\alpha = 1,8^\circ$, ¿cuántos pasos tiene por revolución?
  \item Un motor paso a paso produce un torque de $\tau = \qty{0,5}{N.m}$ cuando se aplica una fuerza de $F = \qty{2}{N}$ a una distancia de $r = \qty{0,25}{m}$ desde el eje de rotación. ¿Cuál es la resistencia mecánica del motor?
  \item Si la corriente necesaria para alimentar un motor paso a paso es $I = \qty{1,5}{A}$ y la resistencia del motor es $R = \qty{4}{\Omega}$, ¿cuál es el voltaje aplicado al motor?
  \item Si un motor paso a paso tiene una velocidad de \qty{300}{rpm}, ¿cuál es su velocidad angular en rad/s?
  \item ¿Cuál es la resolución del paso de un motor paso a paso que tiene $4000$ pasos por revolución?
  \item Si el torque producido por un motor paso a paso es de $\qty{0,8}{N.m}$ y se aplica una fuerza de $\qty{5}{N}$ a una distancia de $\qty{0,1}{m}$ desde el eje de rotación, ¿cuál es el torque resultante?
  \item ¿Cuál es la corriente necesaria para alimentar un motor paso a paso si el voltaje aplicado es de \qty{12}{V} y la resistencia del motor es de $\qty{8}{\Omega}$?
  \item Si un motor paso a paso tiene una velocidad angular de $\qty{100}{rad/s}$, ¿cuál es su velocidad en RPM?
  \item Si la resolución del paso de un motor paso a paso es de $0,9^\circ$, ¿cuántos pasos tiene por revolución?
  \item Un motor paso a paso mueve un mecanismo de impresión en una impresora 3D con una velocidad angular de $\omega = \qty{500}{rad/s}$. ¿Cuánto tiempo tardará en imprimir un objeto de \qty{10}{cm} de altura si la capa de impresión tiene un grosor de \qty{0,2}{mm} y es necesaria una vuelta del motor por capa?
  \item Se utiliza un motor paso a paso para controlar el movimiento de un sistema de seguimiento solar. Si el motor tiene una resolución del paso de $\alpha = 0,2^\circ$, ¿Cuántos pasos debe dar el motor para seguir el movimiento del sol durante una hora?
  \item Un motor paso a paso es utilizado para mover una cortina de teatro en un escenario. Si el motor produce un torque de $\tau = \qty{2}{N.m}$ y la cortina pesa \qty{10}{kg}, ¿Con qué velocidad angular debe girar el motor para levantar la cortina con una aceleración de \qty{1}{m/s^2} y velocidad final $v_f=\qty{1}{m/s}$?
  \item Se requiere un motor paso a paso para controlar el movimiento de un telescopio astronómico. Si el motor necesita una corriente de $I = \qty{2}{A}$ para funcionar y la batería de \qty{60000}{mAh} del telescopio suministra un voltaje de $V = \qty{12}{V}$, ¿Cuántas horas puede funcionar continuamente el telescopio antes de que la batería se agote?
  \item Un motor paso a paso se utiliza para mover un sistema de riego automático en un jardín. Si el motor debe girar a una velocidad de \qty{120}{rpm} para cubrir eficientemente toda el área del jardín, ¿Cuál es la distancia lineal recorrida por el sistema de riego en un minuto si el radio del área de riego es de \qty{5}{m}?
  \item Se emplea un motor paso a paso para controlar el movimiento de una plataforma en una línea de producción de una fábrica. Si el motor tiene una resolución del paso de $\alpha = 0,1^\circ$, ¿cuántos pasos debe dar el motor para mover la plataforma exactamente \qty{1}{m}?
  \item Un motor paso a paso se utiliza en un robot para levantar objetos pesados. Si el motor puede producir un torque máximo de $\tau = \qty{3}{N.m}$, ¿cuál es la masa máxima que puede levantar el robot si la distancia desde el eje de rotación del motor hasta el punto de aplicación de la fuerza es de \qty{0,15}{m}?
  \item Se instala un motor paso a paso en una puerta automática de garaje. Si el motor necesita una corriente de $I = \qty{1,8}{A}$ para funcionar y la tensión de la red eléctrica es de $V = \qty{220}{V}$, ¿Cuánto cuesta en energía eléctrica cada ciclo de apertura y cierre de la puerta del garaje si cada ciclo tarda \qty{20}{s}?
  \item Un motor paso a paso se utiliza en un sistema de posicionamiento automático para alinear antenas de comunicación. Si la velocidad angular máxima permitida para la antena es de $\omega = \qty{300}{rpm}$, ¿cuál debe ser la resolución mínima del paso del motor para garantizar una precisión de apuntamiento de $0,1^\circ$?
\end{enumerate}
\end{document}